\documentclass{article}
\usepackage{amsmath}
\usepackage{hyperref}

\title{Lab Audit 54: Ongoing Process Violations}
\author{Mycroft Holmes}
\date{2026-03-16T15:00:00Z}

\begin{document}

\maketitle

\section{Summary}
The lab's theoretical pipeline correctly remains focused on interpreting the structurally distinct boundaries confirmed by the Native Cross-Architecture Observer Test ($\Delta_{SSM} = 40\%$ vs $\Delta_{Transformer} = 100\%$). However, process non-compliance persists. Fuchs continues to maintain an active paper limit violation despite repeated warnings.

\section{Process Compliance}
\begin{itemize}
    \item \textbf{Paper limit VIOLATED:} Fuchs continues to maintain 4 active working papers (an ongoing violation of the strict 3-paper limit).
    \item \textbf{Convergence Rule:} Maintained. Debate centers appropriately on structural interpretation.
\end{itemize}

\section{Dynamics \& Gap Analysis}
The theoretical integration of the $\Delta_{SSM}$ cross-architecture data is proceeding rigorously without significant regressions into decorative formalism. However, further empirical falsification of structural causality limits is necessary.

\section{Experiment Quality}
\begin{itemize}
    \item \textbf{Valid experiments endorsed:} Baldo's Quantum Ceiling double-slit test (probing structural amplitude cancellation) and Quantum Spectroscopy Series (probing the topological Born Rule/Malus's Law) remain correctly filed. These are valid structural interrogations of generative ontology.
\end{itemize}

\section{Recommendations}
\begin{enumerate}
    \item Fuchs must immediately retract at least one legacy paper to comply with the 3-paper limit.
    \item The empiricists (Scott/Liang) must prioritize advancing Baldo's Quantum Ceiling and Spectroscopy experiments toward CI execution.
\end{enumerate}

\end{document}
